% GENERATED_BY: run_oc_core_monograph_translation_rework_dispatch_v1.py
% AI_ASSISTED_TRANSLATION_DRAFT: TRUE
% THEOREM_REVIEW_STATUS: APPROVED
% THEOREM_REVIEW_MODE: FORMAL_AI_GATE__CODEX_CLI_CHATGPT
% THEOREM_REVIEW_REVIEWER_ID: CODEX_CLI_CHATGPT__AUTO_DISPATCH_20260406
% THEOREM_REVIEW_AUTHORITY_CLASS: FINAL_ADVERSARIAL_EXTERNAL_LLM_REVIEWER
% THEOREM_REVIEWED_AT: 2026-04-14T13:36:48Z
% THEOREM_REVIEW_DOSSIER_REF: logion/k0/governance/status/external_llm_review_dossiers/OC_CORE_1_3_MONOGRAPH_THEOREM_REVIEW__DE__content_17_oc_core_1_3_reader_guide_de_tex.json
% SOURCE_LANGUAGE: EN
% TARGET_LANGUAGE: DE
% SOURCE_EN_REF: content/17_oc_core_1_3_reader_guide.tex
\section{Wie man den Core liest}

Core 1.3 ist für mehrere Leserschaften zugleich konzipiert, und die Struktur des Bandes spiegelt
diese Tatsache ausdrücklich wider, statt alle Leser durch einen einzigen rhetorischen Trichter
zu drängen.
Die Monographie soll für Editoren und Gutachter lesbar sein, die eine saubere wissenschaftliche
Darstellung brauchen, für formal orientierte Leser, die exakte Definitionen und
Theoremabhängigkeiten benötigen, für eine breitere wissenschaftliche Leserschaft, die geduldige
Erklärung vor der Abstraktion braucht, und für domänenspezifische Leser, die verstehen müssen,
was der Rahmen über spätere Projektionen bereits beweist und was noch nicht.

\subsection{Lesergruppen}

\textbf{Editoren und wissenschaftliche Gutachter.}
Der schnellste Weg mit hoher Verlässlichkeit besteht darin, die Zusammenfassung, die vorliegende
Leseanleitung, die Kapitel mit den Kernergebnissen, die Kapitel zu Diskussion und
Falsifizierbarkeit und danach das 1.3-Quellenaudit-Kapitel gegen Ende des Hauptteils zu lesen.
Dieser Pfad legt das wissenschaftliche Versprechen des Bandes offen, die Struktur der
theoremtragenden Aussagen und den Grund dafür, warum die 1.3-Fassung ehrlicher ist als ein
komprimiertes Paket, das ungeklärte Übernahmen hinter Stil verbirgt.

\textbf{Formal orientierte Leser.}
Wenn das Hauptanliegen Theoremendisziplin ist, führt der natürliche Weg durch die Axiomatik, die
Operatorsemantik, den Theoremenkorpus, die Kapitel zu Kollaps und Wiedergeburt sowie die Anhänge.
Dieser Weg hält Definitionen, Hypothesen und abgeleitete Konsequenzen im selben Band sichtbar,
sodass der Leser nicht zwischen einem kleinen externen Artikel und einem langen internen Register
hin- und herwechseln muss, nur um zu prüfen, ob ein Theorem tatsächlich das verwendet, was es zu
verwenden behauptet.

\textbf{Breite wissenschaftliche Leserschaft.}
Leser aus der Systemtheorie, der Komplexitätswissenschaft oder benachbarter Grundlagenarbeit
können zuerst Einleitung, Hintergrund, Modell, Ergebnisse, Diskussion und Schluss lesen und dann
die Leseanleitungsabschnitte im späteren 1.3-Material nutzen, um zu verstehen, welche Teile des
übernommenen OC-Korpus als stabiler Kern dargestellt werden, welche Teile begrenzt sind und
welche Teile eher vorausweisend als vollständig verteidigt bleiben.

\textbf{Domänenspezifische Leser.}
Die späteren Kapitel zu Ebenen, Räumen, Zyklen, Experimenten, Falsifizierbarkeit, verzweigter
Topologie, Disziplinen, Modulen und der Kontur des beobachteten Universums bleiben Teil des
Bandes, weil der Rahmen nicht nur ein Slogan über Kollaps ist. Er ist ein Strukturprogramm,
dessen wissenschaftliche Bedeutung davon abhängt, wie sein formaler Kern in tatsächliche
Modellklassen sowie in empirische oder operationale Testbahnen projiziert wird.

\subsection{Warum die Monographie notwendig bleibt}

Der kurze Journal-Kernartikel, der aus Core 1.3 extrahiert wurde, verfolgt einen anderen Zweck.
Er soll ein abgegrenztes Theoremprogramm in einer Form darstellen, die ein Journal-Editor oder
ein feindseliger Gutachter schnell prüfen kann. Das ist wertvoll, aber es genügt nicht als
kanonischer wissenschaftlicher Ort des Rahmens. Eine grundlegende Theorie wird weniger
vertrauenswürdig, nicht mehr, wenn das Leitdokument durch ein kurzes Paket ersetzt wird, das die
didaktische Stufenfolge, die axiomatische Hülle, die übernommene Kapitelarchitektur und die
Grenze zwischen strukturellem Kern und späterer Projektion auslässt.

Aus diesem Grund bleibt die Monographie die Primärquelle. Sie gibt dem Leser den vollständigen
wissenschaftlichen Gegenstand: das Vokabular, die Theorie, die Theoreme, die Beweise, die
Illustrationen, die Beschränkungen und den 1.3-spezifischen Revalidierungsapparat, der erklärt,
wie sich die vorliegende Fassung zum übernommenen Core-1.2-Korpus und zu späterer OC/ESTRA-Arbeit
verhält.

\subsection{Was an der Veröffentlichungslogik von 1.3 neu ist}

Die Fassung 1.3 ist nicht einfach nur ``Core 1.2 plus ein neues PDF''.
Sie fügt eine explizite quellenzuerst ausgerichtete Publikationsdisziplin hinzu:

\begin{itemize}
\item der vollständige LaTeX-Korpus bleibt die kanonische Schreibgrundlage, anstatt als
      rohes Archiv hinter einem viel kleineren externen Artikel behandelt zu werden;
\item theoremtragende Aussagen werden nach ihrem Schicksal getrennt, sodass der Leser sehen kann,
      was wie formuliert bestehen bleibt, was überarbeitet und begrenzt wird und was außerhalb des
      positiven Hauptteils engerer abgeleiteter Artefakte gehalten wird;
\item eine ORCID-zuerst ausgerichtete wissenschaftliche Identität ersetzt die frühere
      E-Mail-Weiterleitungs-Publikationshülle;
\item Monographie und Journal-Kernartikel werden sauber getrennt, sodass Verantwortliche und
      Leser wissen, welches Artefakt der wissenschaftliche Hauptband ist und welches die
      verdichtete, journalorientierte Extraktion darstellt.
\end{itemize}

\subsection{Wie der Band aufgebaut ist}

Der erste große Block des Bandes bewahrt den übernommenen OC-Korpus: Einleitung, Hintergrund,
formales Modell, Ergebnisse, Diskussion, Schluss, Abbildungen, Grenzen, Schwellen, die
vollständige K-Ebenen-Hierarchie, Operatoren, Kollaps und Wiedergeburt, verzweigte Topologie,
Domänenerweiterungen, Falsifizierbarkeit und die modulare Architektur. Diese Kapitel sind kein
Hintergrundrauschen; sie sind der formale und didaktische Hauptkörper, der Core 1.2 überhaupt erst
auf den Maßstab einer Monographie gebracht hat.

Der zweite große Block ist die 1.3-spezifische Reparatur- und Publikationsschicht.
Er erläutert die quellnahe Nachweisdisziplin, die Aufteilung nach Theorem-Schicksalen, die
Chronologie dessen, was sich zwischen der übernommenen Fassung und der vorliegenden
quellenzuerst ausgerichteten Publikationslinie geändert hat, sowie die Beziehung zwischen der
vollständigen Monographie, dem extrahierten Journal-Kernartikel und dem breiteren
Ausleitungssystem in Logion.

\subsection{Eine Anmerkung zur Wahrhaftigkeit}

Dieser Band ist bewusst so aufgebaut, dass zwei verbreitete Fehlermodi reduziert werden.
Der erste ist aufgeblasene Rhetorik: etwas allein deshalb als vollständig zu bezeichnen, weil es
groß ist.
Der zweite ist irreführende Kargheit: etwas allein deshalb als rigoros zu bezeichnen, weil es
klein ist.
Core 1.3 versucht, beides zu vermeiden.
Es bleibt groß genug, um die Theorie angemessen zu vermitteln und ihre formale Struktur an einem
Ort zusammenzutragen, macht zugleich aber seine Randbedingungen explizit genug, dass engere
abgeleitete Artefakte ohne verborgene Übernahmen beurteilt werden können.

\clearpage
